\documentclass[11pt,a4paper]{article}
\usepackage[margin=1in]{geometry}
\usepackage{booktabs}
\usepackage{graphicx}
\usepackage{amsmath,amssymb}
\usepackage{hyperref}
\usepackage[font=small]{caption}
\usepackage{xcolor}
\usepackage{enumitem}

\setlist[itemize]{topsep=2pt,itemsep=1pt,parsep=0pt}

\title{Grokking Under Federated Averaging:\\Robustness and Failure Modes}
\author{}
\date{}

\begin{document}
\maketitle
\vspace{-2em}

\section{Setup}

\begin{itemize}
  \item \textbf{Task}: modular addition $(a+b) \bmod 97$, training fraction $\alpha{=}0.5$
  \item \textbf{Model}: GrokNet --- 2-layer MLP, no biases, quadratic activation, mean-field init, $N{=}128$, MSE loss
  \item \textbf{Centralised baseline}: full-batch GD, $\mathrm{lr}{=}50$, $\mathrm{wd}{=}0$
  \item \textbf{Federated}: Flower simulation, $K$ clients, 2000 rounds, 5 local epochs (defaults)
  \item \textbf{Partition schemes}: IID, operand-based ($a \bmod K = k$), target-based ($(a{+}b) \bmod K = k$)
\end{itemize}

\section{Experiment 1: Federated vs.\ Centralised Baseline}

\begin{itemize}
  \item FedAvg (IID, $K{=}5$, full participation) reproduces the centralised grokking trajectory step-for-step (Figure~\ref{fig:fed-vs-central})
  \item Both reach 100\% test accuracy at ${\sim}8200$ gradient steps with identical IPR dynamics ($0.02 \to 0.34$)
  \item All three partition strategies (IID, operand, target) produce indistinguishable curves at $K{=}5$
\end{itemize}

\begin{figure}[h]
\centering
\includegraphics[width=0.85\textwidth]{../results/federated/fed_vs_central_iid_K5.png}
\caption{Federated (IID, $K{=}5$) vs.\ centralised training. Trajectories are effectively identical.}
\label{fig:fed-vs-central}
\end{figure}

\section{Experiment 2: Stress-Testing Federated Grokking}

Three axes probed to find where federated grokking breaks.

\subsection{Client Scaling (IID, full participation)}

\begin{table}[h]
\centering
\caption{Client scaling. All configurations grok to 100\%.}
\label{tab:clients}
\begin{tabular}{rrrrr}
\toprule
$K$ & Samples/client & Grok step & Final test acc & Final IPR \\
\midrule
5   & 940 & 8215 & 100.0\% & 0.342 \\
10  & 470 & 8220 & 100.0\% & 0.342 \\
20  & 235 & 8240 & 100.0\% & 0.340 \\
50  & 94  & 8280 & 100.0\% & 0.332 \\
97  & 48  & 8430 & 100.0\% & 0.311 \\
\bottomrule
\end{tabular}
\end{table}

\begin{itemize}
  \item $K{=}97$ (48 samples/client) still groks to 100\%
  \item Grok step increases only 2.6\% from $K{=}5$ to $K{=}97$
  \item Grokking is immune to client count under IID
\end{itemize}

\subsection{Partial Participation (K=20, target partition)}

\begin{table}[h]
\centering
\caption{Partial participation under non-IID target partition ($K{=}20$).}
\label{tab:participation}
\begin{tabular}{rrrrrr}
\toprule
Fraction & Clients/round & Grok step & Final test acc & Final IPR \\
\midrule
1.0 & 20 & 9585  & 98.5\% & 0.214 \\
0.8 & 16 & 9600  & 97.5\% & 0.218 \\
0.6 & 12 & 9725  & 96.2\% & 0.220 \\
0.4 & 8  & 9900  & 93.9\% & 0.222 \\
0.2 & 4  & 9720  & 94.5\% & 0.242 \\
\bottomrule
\end{tabular}
\end{table}

\begin{itemize}
  \item \textbf{Only axis that degrades grokking}
  \item Non-IID target + low participation ($f_t{=}0.2$) caps accuracy at ${\sim}94\%$ vs.\ 100\% under IID
  \item Cause: not all target residue classes represented each round $\Rightarrow$ systematic under-representation
  \item Grokking still occurs --- accuracy jumps from ${\sim}1\%$ to ${\sim}94\%$ at roughly the same step
\end{itemize}

\subsection{Extreme Local Epochs (K=5, IID)}

\begin{table}[h]
\centering
\caption{Extreme local epochs ($K{=}5$, IID, full participation).}
\label{tab:local-epochs}
\begin{tabular}{rrrrr}
\toprule
Local epochs & Rounds & Grok step & Final test acc & Final IPR \\
\midrule
50  & 400 & 8300  & 100.0\% & 0.457 \\
100 & 200 & 8400  & 100.0\% & 0.446 \\
200 & 100 & 9000  & 100.0\% & 0.426 \\
\bottomrule
\end{tabular}
\end{table}

\begin{itemize}
  \item $\mathrm{le}{=}200$ (only 100 communication rounds) still groks to 100\%
  \item Higher local epochs produce \emph{higher} IPR ($0.46$ vs.\ $0.34$ at $\mathrm{le}{=}5$)
  \item Extended local training strengthens periodic Fourier features rather than harming them
\end{itemize}

\begin{figure}[h]
\centering
\includegraphics[width=\textwidth]{../results/federated/fed_breaking_point.png}
\caption{Breaking point: $3{\times}3$ grid. Rows: client scaling, partial participation, extreme local epochs. Columns: test accuracy, IPR, train accuracy.}
\label{fig:breaking}
\end{figure}

\section{Experiment 3: FedProx Sweep}

\begin{itemize}
  \item FedProx adds proximal term $\frac{\mu}{2}\|w - w^{\mathrm{global}}\|^2$ to each client's local objective
  \item Sweep: $\mu \in \{0, 0.001, 0.01, 0.1, 1.0, 10.0\}$, $K{=}5$, IID, full participation
\end{itemize}

\begin{table}[h]
\centering
\caption{FedProx sweep. Sharp transition between $\mu{=}0.001$ (groks) and $\mu{=}0.01$ (fails).}
\label{tab:fedprox}
\begin{tabular}{rlrrrl}
\toprule
$\mu$ & Label & Final test acc & Grok step & $\|W_1\|_F$ & Status \\
\midrule
0     & FedAvg    & 100.0\% & 8215  & 23.80 & Groks \\
0.001 & FedProx   &  99.7\% & 9075  & 22.39 & Groks (10\% delayed) \\
0.01  & FedProx   &   0.1\% & ---   & 15.14 & Weight growth suppressed \\
0.1   & FedProx   &   1.0\% & ---   & NaN   & Diverges \\
1.0   & FedProx   &   1.0\% & ---   & NaN   & Explodes (round 2) \\
10.0  & FedProx   &   1.0\% & ---   & NaN   & Explodes (round 1) \\
\bottomrule
\end{tabular}
\end{table}

\begin{itemize}
  \item \textbf{Sharp phase transition} between $\mu{=}0.001$ and $\mu{=}0.01$
  \item $\mu{=}0.001$: groks with 10\% delay (step 9075 vs.\ 8215), slightly lower weight norm (22.4 vs.\ 23.8)
  \item $\mu{=}0.01$: weight norm capped at 15.1 (vs.\ 23.8), only 44\% train accuracy, no grokking
  \item $\mu \ge 0.1$: proximal gradient $\mu(w - w^{\mathrm{global}})$ $\times$ large lr ($50$) $\Rightarrow$ immediate numerical divergence
\end{itemize}

\begin{figure}[h]
\centering
\includegraphics[width=\textwidth]{../results/federated/fed_fedprox_sweep.png}
\caption{FedProx sweep. Left: test accuracy. Centre/Right: IPR and weight norm ($\mu \le 0.01$ only; $\mu \ge 0.1$ diverges).}
\label{fig:fedprox}
\end{figure}

\section{Discussion}

\paragraph{Grokking is a property of the loss landscape, not the optimisation path.}
\begin{itemize}
  \item FedAvg's noisy, distributed optimisation does not prevent the phase transition
  \item Tested: up to 97 clients, as few as 4 clients/round, up to 200 local epochs
  \item The grokking mechanism (periodic Fourier mode learning via weight norm growth) depends on loss landscape geometry, which is preserved regardless of how gradient steps are distributed
\end{itemize}

\paragraph{The proximal term directly opposes grokking.}
\begin{itemize}
  \item Grokking requires weight norms to grow from ${\sim}11$ to ${\sim}24$ (``lazy'' $\to$ ``rich'' regime)
  \item FedProx's $\frac{\mu}{2}\|w - w^{\mathrm{global}}\|^2$ acts as additional weight decay, anchoring parameters near the global model
  \item Even $\mu{=}0.01$ suppresses growth below the critical threshold ($\|W_1\|_F \approx 15$ vs.\ $24$)
  \item \textbf{Fundamental tension}: the regularisation that makes FedProx effective for heterogeneous FL (reducing client drift) directly conflicts with the weight norm growth that enables grokking
\end{itemize}

\paragraph{Partial participation is the only weak point.}
\begin{itemize}
  \item Only non-IID target partition + low participation ($f_t{=}0.2$, $K{=}20$) degrades accuracy (to ${\sim}94\%$)
  \item Not a failure of grokking itself --- the sharp transition still occurs
  \item Information-theoretic limitation: with 4 of 20 clients/round, some target residue classes are systematically under-represented
\end{itemize}

\end{document}
